\chapter{極限と微分}\label{章:極限と微分}
\section{数列の極限を言葉と式で捉える}\label{節:数列極限}
前章では，長方形の幅を小さくするときに $n$ を大きくし，割線を接線へ近づけるときに $h$ を $0$ に近づけた。どちらも，変化する数がある値に近づく様子を調べている。まず，番号 $n=1,2,3,\ldots$ に従って並んだ数列 $a_1,a_2,a_3,\ldots$ から考えよう。

\begin{definition}{数列の収束}
$n$ を十分大きくすれば，$a_n$ と一定の実数 $\alpha$ との差を，どのように小さく指定した正の数よりも小さくでき，それ以降も小さいままであるとき，数列 $\{a_n\}$ は $\alpha$ に\textbf{収束する}という。
\[
 \lim_{n\to\infty}a_n=\alpha
 \qquad\text{または}\qquad a_n\to\alpha\quad(n\to\infty)
\]
と書き，$\alpha$ を極限値という。
\end{definition}

「それ以降も」という部分は大切である。たとえば $a_n=(-1)^n$ は $1$ を何度も取るが，その間に $-1$ も取るので，$1$ に収束しない。極限は，一部の項だけが近いという意味ではない。また，$a_n=\frac1n$ はどの項も $0$ ではないが，$0$ に収束する。極限値を実際に取る必要もない。

定義を数式で書くと，「任意の $\varepsilon>0$ に対して，ある自然数 $N$ が存在し，$n\ge N$ なら $|a_n-\alpha|<\varepsilon$ となる」ということである。たとえば $a_n=\frac1n$ なら，$N>\frac1\varepsilon$ となる自然数を選べばよい。すると $n\ge N$ に対して
\[
 \left|\frac1n-0\right|=\frac1n\le\frac1N<\varepsilon
\]
が成り立つ。極限の記号は，このような誤差の条件を短く表している。

数列が限りなく大きくなる場合には $a_n\to+\infty$，限りなく小さくなる場合には $a_n\to-\infty$ と書く。ただし，$\infty$ は実数の極限値ではない。このような発散と，$(-1)^n$ のような振動による発散は区別する。

収束する数列については，和，差，積の極限をそれぞれの極限から計算できる。商では分母の極限が $0$ でないことが必要である。具体的には $a_n\to\alpha$，$b_n\to\beta$ のとき，
\[
 a_n+b_n\to\alpha+\beta,\qquad
 a_nb_n\to\alpha\beta,\qquad
 \frac{a_n}{b_n}\to\frac\alpha\beta\quad(\beta\ne0)
\]
となる。商の式は十分大きな $n$ で分母が $0$ でない範囲で使う。これらの極限法則と，次の挟み撃ちの原理は，極限の定義から証明できる基本性質として用いる。

\begin{theorem}{挟み撃ちの原理}
十分大きな $n$ について $a_n\le b_n\le c_n$ が成り立ち，$a_n\to\alpha$，$c_n\to\alpha$ なら，$b_n\to\alpha$ である。
\end{theorem}

実際，$a_n,c_n$ がともに $\alpha$ との差を $\varepsilon$ より小さくすれば，その間の $b_n$ も $\alpha-\varepsilon<b_n<\alpha+\varepsilon$ に入る。前章の長方形による面積の計算は，この原理を使う議論だった。

\begin{example}{極限を求めやすい形に直す}
次の極限を求めよ。
\[
 \text{(1)}\quad\lim_{n\to\infty}\frac{3n^2-n+2}{2n^2+5},
 \qquad
 \text{(2)}\quad\lim_{n\to\infty}\bigl(\sqrt{n^2+3n}-n\bigr).
\]
\end{example}
(1) は分子も分母も大きくなるため，そのままでは商の極限法則を適用できない。最も高い次数 $n^2$ で分子と分母を割ると，
\[
 \frac{3n^2-n+2}{2n^2+5}
 =\cfrac{3-\cfrac1n+\cfrac2{n^2}}{2+\cfrac5{n^2}}
 \longrightarrow\frac32.
\]
分母の極限が $2\ne0$ となったので，ここで商の法則を使える。

(2) は大きい数同士の差である。$\sqrt{n^2+3n}$ と $n$ がともに大きくなることだけでは，差の極限は決まらない。共役な式を掛けて差を積に変えると，
\begin{align*}
 \sqrt{n^2+3n}-n
 &=\frac{3n}{\sqrt{n^2+3n}+n}\\
 &=\cfrac3{\sqrt{1+\cfrac3n}+1}\longrightarrow\frac32.
\end{align*}
ここでは平方根の連続性も使った。大きい数同士の差が，有限の値へ近づく場合がある。

\subsection{単調で有界な数列}
数列が\textbf{上に有界}であるとは，すべての $n$ について $a_n\le M$ を満たす一定の実数 $M$ が存在することをいう。下に有界であることも，同様に一定の下限によって定義する。増加を続けていても，上に有界なら無限に大きくなることはできない。

\begin{theorem}{単調有界数列の収束：採用する事実}
単調増加する数列が上に有界なら，有限の極限値へ収束する。単調減少する数列が下に有界なら，同じく有限の極限値へ収束する。
\end{theorem}

この定理の完全な証明には，実数の連続性に関する性質が必要になる。本書では実数の基礎理論までさかのぼらず，この定理を採用する。ただし，適用するときに単調性と有界性の両方を確かめることは省略しない。収束すると予想した式に極限値を代入するだけでは，収束の証明にはならないからである。

\section{等比数列と無限級数}\label{節:等比級数}
等比数列 $1,r,r^2,\ldots$ の有限和を $S_n=1+r+\cdots+r^{n-1}$ とおく。$rS_n=r+r^2+\cdots+r^n$ を引けば，中間の項が消えるため，
\[
 (1-r)S_n=1-r^n,
 \qquad
 S_n=\frac{1-r^n}{1-r}\quad(r\ne1)
\]
を得る。まず有限の項を足した式を確定させ，その後で $n$ を大きくするのが順序である。

$|r|<1$ なら $r^n\to0$ なので，
\[
 \sum_{k=0}^{\infty}r^k
 :=\lim_{n\to\infty}\sum_{k=0}^{n-1}r^k
 =\frac1{1-r}
\]
となる。無限個の項を一度に足しているのではなく，\textbf{部分和の数列の極限}を定義している。一般に $\sum_{k=1}^{\infty}a_k$ が収束するとは，部分和 $\sum_{k=1}^{n}a_k$ が有限の極限値をもつことをいう。

なぜ $|r|<1$ で $r^n\to0$ になるのかも確認できる。$0<q=|r|<1$ とし，$\frac1q=1+c$，$c>0$ とおく。二項定理から $(1+c)^n\ge1+nc$ なので，
\[
 0\le |r^n|=q^n\le\frac1{1+nc}\longrightarrow0.
\]
$r=0$ の場合は直接明らかである。

$r=1$ なら部分和は $n$ となり，$r=-1$ なら部分和は $1,0,1,0,\ldots$ と振動する。$|r|>1$ では各項 $r^n$ さえ $0$ に近づかないため，級数は収束しない。級数が収束するなら，隣り合う部分和の差である一般項は必ず $0$ に近づく。この条件は必要条件であり，一般の級数では十分条件ではない。たとえば調和級数の発散は\chapterref{章:積分と数列}で積分との比較から調べる。

\begin{example}{無限和の値と打切り誤差}
$\sum_{k=0}^{\infty}\left(-\frac13\right)^k$ の和を求めよ。また，$k=0$ から $n-1$ までの和で近似するとき，誤差の絶対値を $10^{-4}$ 未満にする自然数 $n$ を一つ挙げよ。
\end{example}
公比の絶対値は $\frac13<1$ だから，和は
\[
 S=\cfrac{1}{1+\cfrac13}=\frac34.
\]
有限和と極限値との差は，有限和の公式から
\[
 S-S_n=\cfrac{\left(-\cfrac13\right)^n}{1+\cfrac13},
 \qquad
 |S-S_n|=\frac{3}{4\cdot3^n}.
\]
$n=9$ では $\frac3{4\cdot19683}<10^{-4}$ となるので，$9$ 項まで足せばよい。交互に符号が変わるため誤差の符号も変わるが，誤差の大きさはこの一つの式で評価できる。

\section{関数の極限と連続性}\label{節:関数極限連続}
数列では自然数 $n$ を大きくした。関数では，実数 $x$ をある点 $a$ に近づける。$x\ne a$ のまま $x$ を $a$ に十分近づければ，$f(x)$ を $L$ にどれだけでも近づけられるとき，
\[
 \lim_{x\to a}f(x)=L
\]
と書く。正確には，任意の $\varepsilon>0$ に対して，ある $\delta>0$ を選び，$0<|x-a|<\delta$ なら $|f(x)-L|<\varepsilon$ となることである。関数の定義域に属する $x$ だけを考える。

$x>a$ の側から近づける右極限を $\lim_{x\to a+}f(x)$，$x<a$ の側から近づける左極限を $\lim_{x\to a-}f(x)$ と書く。両側で近づけられる点では，左右の極限が同じ値に一致して初めて両側の極限が存在する。たとえば $\frac{|x|}{x}$ は $x\to0+$ で $1$，$x\to0-$ で $-1$ だから，$x\to0$ の極限は存在しない。

\begin{definition}{連続性}
$f$ が $x=a$ で定義され，$\lim_{x\to a}f(x)=f(a)$ となるとき，$f$ は $a$ で\textbf{連続}であるという。閉区間の端点では，区間の内側からの片側極限を使う。
\end{definition}

極限は $x=a$ 以外の近くの値から決まるのに対し，連続性はその極限と実際の値 $f(a)$ が一致するかどうかも問う。多項式，分母が $0$ でない有理関数，および三角関数・指数関数・対数関数は，それぞれの定義域で連続である。和・積・商・合成による連続性の保存も基本性質として用いる。

\begin{example}{値を補って連続にする}
$x\ne2$ に対して $f(x)=\frac{x^2-4}{x-2}$ と定める。$f(2)$ をどのように定めれば，$x=2$ でも連続になるか。
\end{example}
$x\ne2$ の範囲では約分でき，
\[
 f(x)=\frac{(x-2)(x+2)}{x-2}=x+2.
\]
したがって $\lim_{x\to2}f(x)=4$ である。$f(2)=4$ と定めれば極限と関数値が一致し，連続になる。約分後の $x+2$ は $x=2$ でも値をもつが，元の分数はそのままでは定義されない。式が等しい範囲と，関数の定義域を分けて考える必要がある。

後で微分の定理を証明するために，連続関数について次の二つの存在定理を使う。

\begin{theorem}{連続関数の存在定理：採用する事実}\label{定理:連続関数存在}
$f$ が閉区間 $[a,b]$ 上で連続なら，次が成り立つ。
\begin{enumerate}[label=(\arabic*),leftmargin=*]
 \item $f$ は区間内で最大値と最小値をとる。
 \item $f(a)$ と $f(b)$ の間にある任意の値 $k$ に対して，$f(c)=k$ となる $c\in[a,b]$ が存在する。これを\textbf{中間値の定理}という。
\end{enumerate}
\end{theorem}

これらの証明も実数の連続性に基づくため，本書では採用する事実とする。条件を外すと結論は崩れる。たとえば $f(x)=x$ は開区間 $(0,1)$ では最大値も最小値もとらない。端の値 $0,1$ に近づくことはできるが，区間の中には到達する点がない。閉区間という条件は，単なる表記上の飾りではない。

\section{微分係数と導関数}\label{節:微分定義}
\begin{definition}{微分係数と導関数}
$f$ が $a$ の近くで定義され，有限の極限
\[
 f'(a)=\lim_{h\to0}\frac{f(a+h)-f(a)}h
\]
が存在するとき，$f$ は $a$ で\textbf{微分可能}であるといい，$f'(a)$ を $a$ における微分係数という。各点の微分係数を対応させた関数 $f'(x)$ を\textbf{導関数}という。
\end{definition}

微分係数は一つの点での数，導関数は点を動かしたときの関数である。$y=f(x)$ のとき $f'(x)$ を $y'$ や $\frac{\dd y}{\dd x}$ とも書く。$\frac{\dd y}{\dd x}$ は通常の分数そのものではないが，「何を何で微分しているか」を明示する記号である。

$f(x)=x^2$ なら，
\[
 \frac{f(a+h)-f(a)}h=2a+h\longrightarrow2a
\]
なので，$f'(a)=2a$，導関数は $f'(x)=2x$ である。点 $(a,f(a))$ における接線は，傾き $f'(a)$ をもつ直線として
\[
 y-f(a)=f'(a)(x-a)
\]
で表される。

微分可能なら連続である。実際，$h\ne0$ に対して
\[
 f(a+h)-f(a)
 =h\frac{f(a+h)-f(a)}h
\]
であり，右辺は $h\to0$ のとき $0\cdot f'(a)=0$ に近づく。一方，連続だから微分可能とは限らない。$f(x)=|x|$ は $0$ で連続だが，差商 $\frac{|h|}{h}$ の右極限は $1$，左極限は $-1$ なので，$0$ で微分できない。折れ曲がった点に一つの接線の傾きを与えようとしても，左右の候補が一致しないのである。

\subsection{微分を一次式で表す}
微分可能性は，
\[
 f(a+h)=f(a)+f'(a)h+h\varepsilon(h),
 \qquad \varepsilon(h)\to0\quad(h\to0)
\]
と書き直せる。実際，$h\ne0$ に対して $\varepsilon(h)=\frac{f(a+h)-f(a)}h-f'(a)$ とおけばよい。最後の $h\varepsilon(h)$ は，$h$ で割ると $0$ に近づく誤差である。「小さい変化に対して関数を接線で近似できる」という意味が，ここに数式で表されている。この形は合成関数の微分を証明するときにも役立つ。

\section{積・商・合成関数の微分}\label{節:微分法則}
基本的な関数を一つずつ定義から微分しても，複雑な式を毎回展開するのは大変である。そこで，式の作られ方と微分の関係を先に調べる。

\begin{theorem}{和・積・商の微分}
$f,g$ が微分可能で，$c$ が定数なら，
\begin{align*}
 (f+g)'&=f'+g',& (cf)'&=cf',\\
 (fg)'&=f'g+fg',&
 \left(\frac fg\right)'&=\frac{f'g-fg'}{g^2}\quad(g\ne0)
\end{align*}
が成り立つ。
\end{theorem}

和と定数倍は差商を分ければ分かる。積の公式では，二つの因子が同時に変わることに注意する。途中の項を一つ足して引けば，
\begin{align*}
 &f(x+h)g(x+h)-f(x)g(x)\\
 &\quad=\{f(x+h)-f(x)\}g(x+h)
       +f(x)\{g(x+h)-g(x)\}.
\end{align*}
これを $h$ で割って極限をとる。微分可能な $g$ は連続だから $g(x+h)\to g(x)$ となり，積の公式が得られる。積を微分して単に $f'g'$ としてはいけない。

逆数については，
\[
 \cfrac{\cfrac1{g(x+h)}-\cfrac1{g(x)}}h
 =-\frac{g(x+h)-g(x)}{h\,g(x+h)g(x)}
 \longrightarrow-\frac{g'(x)}{g(x)^2}.
\]
$g(x)\ne0$ と連続性により，十分近くでは $g(x+h)\ne0$ となる。この逆数の公式と積の公式から，商の公式も導ける。

\begin{theorem}{合成関数の微分}\label{定理:合成微分}
$u=g(x)$ が $x$ で微分可能で，$f$ が $u=g(x)$ で微分可能なら，
\[
 \frac{\dd}{\dd x}f(g(x))=f'(g(x))g'(x)
\]
である。
\end{theorem}

外側の関数の変化率に，内側の変数の変化率を掛ける公式である。証明のため，$\Delta u=g(x+h)-g(x)$ とおく。$g$ の連続性から $\Delta u\to0$ であり，$f$ の微分可能性から
\[
 f(u+\Delta u)-f(u)=f'(u)\Delta u+\Delta u\,\eta(\Delta u),
 \qquad \eta(\Delta u)\to0
\]
と書ける。$\eta(0)=0$ と定めれば，$\Delta u=0$ の場合もこの等式は使える。$h$ で割ると
\[
 \frac{f(g(x+h))-f(g(x))}h
 =\{f'(g(x))+\eta(\Delta u)\}
   \frac{g(x+h)-g(x)}h.
\]
右辺の極限は $f'(g(x))g'(x)$ になる。途中で $\Delta u$ を分母に置いていないので，内側の関数が同じ値を取る場合も含めて証明できた。

\section{基本関数を微分する}\label{節:基本微分}
\subsection{べき関数}
正の整数 $n$ について二項定理を使うと，
\[
 \frac{(x+h)^n-x^n}{h}
 =nx^{n-1}+\binom n2x^{n-2}h+\cdots+h^{n-1}.
\]
$h\to0$ とすれば，
\[
 \frac{\dd}{\dd x}x^n=nx^{n-1}
\]
を得る。定数の微分は $0$ であり，負の整数べきは逆数の微分から同じ形になる。実数べきへの拡張は，対数関数を微分した後で説明する。

\subsection{三角関数と弧度法}
三角関数を微分するとき，角は必ず弧度法で表す。基本になる極限は
\[
 \lim_{h\to0}\frac{\sin h}{h}=1
\]
である。$0<h<\frac\pi2$ では，単位円における三角形，扇形，外側の三角形の面積比較から
\[
 \sin h<h<\tan h
\]
を得る。したがって $\cos h<\frac{\sin h}{h}<1$ であり，挟み撃ちの原理から右極限は $1$ となる。$\frac{\sin h}{h}$ は $h$ の符号を変えても同じだから，左極限も $1$ である。また，
\[
 \frac{\cos h-1}{h}
 =-\cfrac{2\sin^2\cfrac h2}{h}
 =-\sin\frac h2\,\cfrac{\sin\cfrac h2}{\cfrac h2}
 \longrightarrow0.
\]
加法定理に代入すれば，
\begin{align*}
 \frac{\sin(x+h)-\sin x}{h}
 &=\sin x\frac{\cos h-1}{h}
   +\cos x\frac{\sin h}{h}\\
 &\longrightarrow\cos x.
\end{align*}
同様に余弦を微分でき，商の公式から正接も求まる。
\[
 (\sin x)'=\cos x,\qquad
 (\cos x)'=-\sin x,\qquad
 (\tan x)'=\frac1{\cos^2x}\quad(\cos x\ne0).
\]
度数法で $x$ を表すと，弧度への換算係数が微分に現れる。上の簡潔な公式が成り立つのは，弧度法を使っているからである。

\subsection{指数関数と対数関数}
指数関数を微分するためには，底を特別な値 $\e$ に選ぶ。
\begin{theorem}{自然対数の底と基本極限：採用する事実}
数列 $\left(1+\frac1n\right)^n$ は極限 $\e$ をもち，$2<\e<3$ である。この底に対する指数関数は
\[
 \lim_{h\to0}\frac{\e^h-1}{h}=1
\]
を満たす。$\e$ を底とする対数を自然対数といい，本書では $\log x$ と書く。
\end{theorem}

この極限を実数指数の定義から構成する議論は，本書の証明範囲から外す。以後の微分公式は，この基本極限と指数法則から導く。$\e^{x+h}=\e^x\e^h$ なので，
\[
 \frac{\e^{x+h}-\e^x}{h}
 =\e^x\frac{\e^h-1}{h}\longrightarrow\e^x.
\]
よって $(\e^x)'=\e^x$ である。

対数では $x>0$ とし，$k=\log(x+h)-\log x$ とおく。対数の連続性から $h\to0$ のとき $k\to0$ である。また $x+h=x\e^k$ だから，$h=x(\e^k-1)$ となる。したがって
\[
 \frac{\log(x+h)-\log x}{h}
 =\frac1x\frac{k}{\e^k-1}\longrightarrow\frac1x.
\]
以上と合成関数の微分から，
\begin{align*}
 (\e^{u(x)})'&=\e^{u(x)}u'(x),\\
 (\log u(x))'&=\frac{u'(x)}{u(x)}\quad(u(x)>0),\\
 (a^x)'&=a^x\log a\quad(a>0)
\end{align*}
を得る。最後は $a^x=\e^{x\log a}$ を使った。さらに $u(x)\ne0$ の区間では，$u$ の符号を分けると
\[
 (\log|u(x)|)'=\frac{u'(x)}{u(x)}
\]
も成り立つ。絶対値を付けることで，負の値の区間でも実数の対数として定義できる。

$x>0$ では実数 $\alpha$ に対して $x^\alpha=\e^{\alpha\log x}$ と書けるため，
\[
 (x^\alpha)'=\e^{\alpha\log x}\frac\alpha x
 =\alpha x^{\alpha-1}
\]
である。整数べきなら，元の関数が定義される負の $x$ にも既に同じ公式が成り立つ。非整数べきの場合は，まず $x>0$ を基本の定義域として考える。

\begin{example}{式の構造に合わせて微分する}
次の関数を微分し，(1) については $x=2$ における接線も求めよ。
\[
 \text{(1)}\quad f(x)=\frac{x^2+1}{x-1},\qquad
 \text{(2)}\quad g(x)=\e^{x^2}\sin(3x),\qquad
 \text{(3)}\quad y=x^x\quad(x>0).
\]
\end{example}
(1) に商の公式を使うと，
\[
 f'(x)=\frac{2x(x-1)-(x^2+1)}{(x-1)^2}
      =\frac{x^2-2x-1}{(x-1)^2}.
\]
$f(2)=5$，$f'(2)=-1$ なので，接線は $y-5=-(x-2)$，すなわち $y=-x+7$ である。別の計算として，多項式の割り算から
\[
 f(x)=x+1+\frac2{x-1},\qquad
 f'(x)=1-\frac2{(x-1)^2}
\]
としてもよい。微分する前に式を整理すると，計算が短くなる場合がある。

(2) は二つの関数の積であり，それぞれの中に合成がある。まず外側の積の公式を使い，次に内側を微分する。
\[
 g'(x)=2x\e^{x^2}\sin(3x)+3\e^{x^2}\cos(3x).
\]
$2x$ と $3$ は，それぞれ $x^2$ と $3x$ を微分した因子である。

(3) は指数も底も変わるので，定数指数の公式をそのまま使えない。$y>0$ だから対数をとり，$\log y=x\log x$ とすると，
\[
 \frac{y'}y=\log x+1,
 \qquad y'=x^x(\log x+1).
\]
このように対数をとって積やべきを分けてから微分する方法を\textbf{対数微分法}という。対数をとれる条件 $x>0$，$y>0$ を先に確かめている点も大切である。

\section{平均値の定理が局所と区間を結ぶ}\label{節:平均値定理}
微分係数は一点での変化率である。一方，関数が区間全体で増えるかどうかを知るには，離れた二点の値を比べる必要がある。この二つをつなぐのが平均値の定理である。証明の準備として，極値で何が起こるかを確かめよう。

点 $c$ の十分近くで $f(x)\le f(c)$ なら，$f$ は $c$ で\textbf{極大値}をとるという。不等号を逆にしたものが\textbf{極小値}である。極大・極小は近くの点との比較であり，区間全体での最大・最小とは区別する。

$c$ が区間の内点で，$f$ がそこで微分可能かつ極大なら，小さい $h$ に対して $f(c+h)-f(c)\le0$ である。したがって差商は $h>0$ で $0$ 以下，$h<0$ で $0$ 以上となる。左右の極限が一致するには，$f'(c)=0$ でなければならない。極小の場合も同じである。ただし，端点では左右両側を比較できず，この結論をそのまま使えない。

\begin{theorem}{ロルの定理}\label{定理:ロル}
$a<b$ とする。$f$ が閉区間 $[a,b]$ で連続，開区間 $(a,b)$ で微分可能で，$f(a)=f(b)$ なら，
\[
 f'(c)=0
\]
を満たす $c\in(a,b)$ が少なくとも一つ存在する。
\end{theorem}

定理\ref{定理:連続関数存在}により，$f$ は最大値と最小値をとる。$f$ が定数なら，どの内点でも導関数は $0$ である。定数でなければ，最大値または最小値の少なくとも一方は共通の端点値 $f(a)=f(b)$ と異なる。その値をとる点は内点なので，先ほどの議論から，その点で $f'(c)=0$ となる。これで証明できた。

\begin{theorem}{平均値の定理}\label{定理:平均値}
$a<b$ とする。$f$ が閉区間 $[a,b]$ で連続，開区間 $(a,b)$ で微分可能なら，
\[
 f'(c)=\frac{f(b)-f(a)}{b-a}
\]
を満たす $c\in(a,b)$ が少なくとも一つ存在する。
\end{theorem}

右辺は二つの端点を結ぶ割線の傾きである。そこで，その傾きの直線を $f$ から引き，両端の高さをそろえる。
\[
 g(x)=f(x)-\frac{f(b)-f(a)}{b-a}(x-a)
\]
とおけば，$g(a)=g(b)=f(a)$ である。また，$g$ は閉区間で連続，内側で微分可能だから，ロルの定理を使える。ある内点 $c$ で
\[
 0=g'(c)=f'(c)-\frac{f(b)-f(a)}{b-a}
\]
となり，求める等式を得る。

\begin{figure}[htbp]
\centering
\begin{tikzpicture}[x=2.6cm,y=1.1cm,>=Stealth]
 \draw[->] (-0.15,0)--(2.35,0) node[right]{$x$};
 \draw[->] (0,-0.4)--(0,4.7) node[above]{$y$};
 \draw[very thick,domain=0:2.12,samples=70] plot (\x,{\x*\x});
 \draw[blue,thick] (0,0)--(2,4);
 \draw[red,thick,domain=0.4:1.8] plot (\x,{2*\x-1});
 \fill (0,0) circle(1.6pt);
 \fill (2,4) circle(1.6pt) node[above left]{$(2,4)$};
 \fill (1,1) circle(1.6pt) node[below right]{$(1,1)$};
 \draw[dashed] (1,0) node[below]{$1$}--(1,1);
 \draw[dashed] (2,0) node[below]{$2$}--(2,4);
 \node[blue] at (0.63,2.05) {割線};
 \node[red] at (1.8,2.2) {接線};
 \node at (1.35,4.55) {$y=x^2$};
\end{tikzpicture}
\caption{平均変化率に等しい微分係数をもつ点が，区間の内側に存在する。}
\label{図:平均値定理}
\end{figure}

図\ref{図:平均値定理}では，$f(x)=x^2$ の区間 $[0,2]$ における平均変化率は $2$ であり，$f'(1)=2$ となる。一般には，その点が区間の中央とは限らず，一つに決まるとも限らない。

条件の役割も確認しておこう。$f(x)=|x|$ を $[-1,1]$ で考えると，両端の値は等しいが，微分可能な点の導関数は $-1$ または $1$ である。$0$ では微分できないため，ロルの定理も平均値の定理も適用できない。導関数の連続性は定理の仮定に含まれていないが，区間内のすべての点で微分可能であることは必要である。

\subsection{導関数の符号から増減が分かる理由}
区間内の任意の二点 $x_1<x_2$ に平均値の定理を使うと，その間のある $c$ に対して
\[
 f(x_2)-f(x_1)=f'(c)(x_2-x_1)
\]
となる。$x_2-x_1>0$ だから，区間内で $f'>0$ なら $f(x_2)>f(x_1)$ であり，$f$ は狭義に増加する。$f'<0$ なら狭義に減少する。$f'\ge0$ の場合には $f(x_2)\ge f(x_1)$ までが保証され，値が一定の部分もあり得る。

同じ議論から，次の結果も得られる。
\begin{theorem}{導関数が同じ関数の差}\label{定理:導関数同じ}
一つの区間のすべての内点で $F'=G'$ が成り立ち，端点を含む場合にはそこで両関数が連続なら，その区間で $F-G$ は定数である。特に，導関数が恒等的に $0$ なら，元の関数は定数である。
\end{theorem}

実際，$H=F-G$ とおくと $H'=0$ であり，任意の二点について $H(x_2)-H(x_1)=0$ となる。後の積分では，微分して同じ関数になる候補が，定数の違いを除いて同じであることをこの定理で確かめる。「一つの区間」という条件にも意味がある。たとえば $x<0$ で $H(x)=0$，$x>0$ で $H(x)=1$ と定めれば，定義域のすべての点で $H'=0$ だが，左右の区間に同じ定数を使うことはできない。

\section{増減・凹凸からグラフを読む}\label{節:増減凹凸}
$f'(c)=0$ は極値の候補を探す条件であり，極値を保証する条件ではない。$f(x)=x^3$ は $f'(0)=0$ だが，$0$ の前後でも増加し続ける。極大・極小を判定するには，その点の前後で導関数の符号がどう変わるかを調べる。$+$ から $-$ に変われば極大，$-$ から $+$ に変われば極小である。

導関数をもう一度微分した $f''=(f')'$ を\textbf{第二次導関数}という。$f''>0$ なら接線の傾き $f'$ が増加し，グラフは\textbf{下に凸}になる。$f''<0$ なら傾きが減少し，グラフは\textbf{上に凸}になる。たとえば $y=x^2$ は下に凸で，$y=-x^2$ は上に凸である。

この呼び方を，接線との位置関係から確かめられる。$f''>0$ の区間で $a$ を固定し，$x>a$ とすると，平均値の定理により
\[
 f(x)-f(a)=f'(c)(x-a)>f'(a)(x-a)
 \qquad(a<c<x)
\]
となる。$x<a$ の場合も，$x<c<a$ で $f'(c)<f'(a)$ となることと $x-a<0$ を合わせれば，同じ結論を得る。したがって，曲線は接線の上側にある。$f''<0$ では不等号が逆になる。

グラフの凹凸が切り替わる点を\textbf{変曲点}という。$f''$ が連続な場合，凹凸が切り替わる場所は $f''=0$ となる点から探せるが，$f''=0$ だけでは十分でない。$y=x^4$ は $x=0$ で $f''=0$ になっても，その前後でともに下に凸である。

\begin{example}{極値と区間全体の最大・最小}
$f(x)=x^3-3x$ の増減と凹凸を調べ，区間 $[-2,3]$ における最大値・最小値を求めよ。
\end{example}
導関数を因数分解すると，符号を調べやすい。
\[
 f'(x)=3x^2-3=3(x-1)(x+1),\qquad f''(x)=6x.
\]
したがって $x<-1$ で増加，$-1<x<1$ で減少，$x>1$ で増加する。$f(-1)=2$ が極大値，$f(1)=-2$ が極小値である。また，$x<0$ で上に凸，$x>0$ で下に凸だから，原点 $(0,0)$ は変曲点である。
\[
\begin{array}{c|ccccccc}
 x&-2&&-1&&1&&3\\ \hline
 f'(x)&&+&0&-&0&+&\\
 f(x)&-2&\nearrow&2&\searrow&-2&\nearrow&18
\end{array}
\]
閉区間での最大・最小は，極値だけを比較しては決まらない。両端も候補になるため，$f(-2)=-2$，$f(3)=18$ を加えて比較する。最大値は $18$ で $x=3$ のとき，最小値は $-2$ で $x=-2,1$ のときにとる。

\begin{figure}[htbp]
\centering
\begin{tikzpicture}[x=2.1cm,y=0.8cm,>=Stealth]
 \draw[->] (-2.3,0)--(2.35,0) node[right]{$x$};
 \draw[->] (0,-3.4)--(0,3.5) node[above]{$y$};
 \draw[very thick,domain=-2.1:2.1,samples=120] plot (\x,{\x^3-3*\x});
 \fill (-1,2) circle(1.5pt) node[above left]{$(-1,2)$};
 \fill (1,-2) circle(1.5pt) node[below right]{$(1,-2)$};
 \fill (0,0) circle(1.5pt) node[above right]{変曲点};
 \draw[dashed] (-1,0) node[below]{$-1$}--(-1,2)--(0,2);
 \draw[dashed] (1,0) node[above]{$1$}--(1,-2)--(0,-2);
 \node at (1.1,2.85) {$y=x^3-3x$};
\end{tikzpicture}
\caption{極値は増減の切り替わり，変曲点は凹凸の切り替わりを表す。}
\label{図:三次増減}
\end{figure}

図\ref{図:三次増減}は原点付近を描いたものである。見えている山の高さ $2$ が，指定区間全体の最大値とは限らないことに注意しよう。最大・最小を求めるときは，まず定義域と対象の区間を確かめ，内点の極値候補，微分できない点，含まれる端点を比較する。連続関数でも，区間が閉じていなければ最大・最小が存在しない場合がある。

\section{方程式と不等式に微分を使う}\label{節:微分方程式不等式}
方程式の根を式で表せなくても，根が存在することや，根の個数は分かることがある。ここでは，連続性が存在を保証し，増減が個数を制限する。不等式では，左右の差を一つの関数にして，その最小値を調べる。

\begin{example}{解の存在と不等式の等号条件}
\begin{enumerate}[label=(\arabic*),leftmargin=*]
 \item 方程式 $x=\cos x$ が区間 $[0,1]$ にただ一つの解をもつことを示せ。
 \item $x>0$ のとき $\log x\le x-1$ を示し，等号成立条件を求めよ。
\end{enumerate}
\end{example}
(1) では $h(x)=x-\cos x$ とおく。$h$ は $[0,1]$ で連続で，
\[
 h(0)=-1<0,\qquad h(1)=1-\cos1>0
\]
である。中間値の定理から，$h(x)=0$ となる点が $(0,1)$ に存在する。一方，この区間では $h'(x)=1+\sin x>0$ だから，$h$ は狭義に増加する。二つの異なる点でともに $0$ をとることはできず，解はただ一つである。グラフの交わりを予想しただけではなく，存在と一意性を別々に示した。

(2) は右辺から左辺を引き，$g(x)=x-1-\log x$ とおく。
\[
 g'(x)=1-\frac1x=\frac{x-1}{x}.
\]
定義域では $x>0$ なので，符号は分子 $x-1$ で決まる。$g$ は $0<x<1$ で減少し，$x>1$ で増加する。したがって，最小値は $g(1)=0$ であり，$g(x)\ge0$ を得る。両側で狭義に増減しているため，等号は $x=1$ のときに限る。

\section{接線で値を近似する}\label{節:一次近似}
\sectionref{節:微分定義}で得た式
\[
 f(a+h)=f(a)+f'(a)h+h\varepsilon(h),\qquad \varepsilon(h)\to0
\]
を使えば，$h$ が小さいとき
\[
 f(a+h)\approx f(a)+f'(a)h
\]
と近似できる。これを\textbf{一次近似}という。右辺は $a$ における接線の値である。計算しやすい点 $a$ を近くに選び，そこからの変化を傾きで見積もる。

\begin{example}{平方根の近似と誤差の向き}
$\sqrt{4.04}$ を一次近似で求め，近似値が実際の値より大きいか小さいかを確かめよ。
\end{example}
$f(x)=\sqrt{x}$ とすると，$f(4)=2$，$f'(4)=\frac1{2\sqrt4}=\frac14$ である。$4.04=4+0.04$ なので，
\[
 \sqrt{4.04}\approx2+\frac14\cdot0.04=2.01.
\]
$f''(x)=-\cfrac1{4x^{\cfrac32}}<0$ だから，グラフは上に凸で，接線による値は実際の値より大きい。数値によっても，$2.01^2=4.0401>4.04$ と確かめられる。

誤差の大きさは，有理化によって直接評価できる。
\[
 2.01-\sqrt{4.04}
 =\frac{2.01^2-4.04}{2.01+\sqrt{4.04}}
 =\frac{0.0001}{2.01+\sqrt{4.04}}.
\]
分母は $4$ より大きいから，$0<2.01-\sqrt{4.04}<0.000025$ である。一次近似の式だけでは，指定した小数桁までの正しさが自動的に保証されるわけではない。この例では別の式で誤差を評価した。より一般的な誤差の扱いは\chapterref{章:近似と広義積分}で学ぶ。

\section{逆関数・媒介変数・陰関数を微分する}\label{節:逆媒介陰微分}
\subsection{逆関数では傾きが逆数になる}
一つの区間で連続かつ狭義単調な関数 $y=f(x)$ は，その値域で逆関数 $x=g(y)$ をもつ。このとき逆関数も連続であるという基本性質を用いる。その区間の内点 $a$ について $b=f(a)$，$f'(a)\ne0$ とし，$y=b+k$ に対応する $x$ を $a+h$ と書けば，$k\to0$ のとき $h\to0$ であり，$k=f(a+h)-f(a)$ である。したがって
\[
 \frac{g(b+k)-g(b)}k
 =\frac{h}{f(a+h)-f(a)}
 =\cfrac{1}{\cfrac{f(a+h)-f(a)}h}
 \longrightarrow\frac1{f'(a)}.
\]
よって
\[
 g'(b)=\frac1{f'(a)},\qquad
 g'(y)=\frac1{f'(g(y))}
\]
となる。分母が $0$ でないことが必要である。たとえば $f(x)=x^2$ を $x>0$ に限ると，逆関数は $g(y)=\sqrt y$ であり，$g'(y)=\frac1{2\sqrt y}$ となる。$x^2$ を実数全体で考えると，正負の二つの $x$ が同じ値に対応して逆関数が定まらないので，先にどちらの枝を使うかを決めなければならない。

\subsection{媒介変数で表した曲線}
曲線が $x=x(t)$，$y=y(t)$ と表されるとき，$t$ を\textbf{媒介変数}という。両関数が連続な導関数をもち，考えている点で $x'(t)\ne0$ なら，その近くで $x(t)$ は狭義単調であり，$t$ を $x$ の関数として表せる。合成関数と逆関数の微分から
\[
 \frac{\dd y}{\dd x}
 =\frac{y'(t)}{x'(t)}
\]
を得る。$x'(t)=0$ の点ではこの商は使えず，垂直接線などを別に調べる必要がある。さらに二階まで微分できる場合は，$x$ での微分が $\frac1{x'(t)}\frac{\dd}{\dd t}$ に対応するので，
\[
 \frac{\dd^2y}{\dd x^2}
 =\frac1{x'(t)}\frac{\dd}{\dd t}
   \left(\frac{y'(t)}{x'(t)}\right)
\]
となる。$\frac{y''(t)}{x''(t)}$ ではない。

\subsection{関係式を保ったまま微分する}
$x,y$ の関係が一つの式で与えられ，$y$ を $x$ の関数として明示していない場合がある。その式が実際に微分可能な関数 $y=y(x)$ を定めている範囲では，両辺を $x$ で微分できる。これを\textbf{陰関数の微分}という。$y$ も $x$ とともに変わるので，たとえば $(y^2)'=2yy'$，$(xy)'=y+xy'$ となる。

\begin{example}{二通りの表し方から接線を求める}
\begin{enumerate}[label=(\arabic*),leftmargin=*]
 \item $x=t^2+1$，$y=t^3-t$ で表される曲線の $t=1$ における接線を求めよ。
 \item 円 $x^2+y^2=25$ の点 $(3,4)$ における接線を求めよ。
\end{enumerate}
\end{example}
(1) では $t=1$ のとき $(x,y)=(2,0)$ である。また，
\[
 \frac{\dd x}{\dd t}=2t,\qquad
 \frac{\dd y}{\dd t}=3t^2-1,\qquad
 \left.\frac{\dd y}{\dd x}\right|_{t=1}
 =\frac{3-1}{2}=1.
\]
したがって接線は $y=x-2$ である。$t$ を消去しなくても，それぞれの変化率の比で傾きを求められた。

(2) では，点 $(3,4)$ の近くで上半円 $y=\sqrt{25-x^2}$ を考えれば，$y$ が微分可能な関数であることを確かめられる。その上で，元の関係式を微分すると
\[
 2x+2yy'=0,\qquad y'=-\frac xy\quad(y\ne0).
\]
点 $(3,4)$ での傾きは $-\frac34$ だから，接線は
\[
 y-4=-\frac34(x-3),\qquad 3x+4y=25
\]
である。円全体を一つの関数 $y=f(x)$ とみなしたわけではなく，指定点の近くの枝を使っている。$y=0$ の点では上の商は使えず，実際に接線は垂直になる。

これで，一点での変化率から，区間の増減や接線の近似までを調べる道具がそろった。次章では，細かい区間での寄与を足して定積分を作り，この微分の議論とつなげる。特に，平均値の定理から得た「導関数が同じなら差は定数」という結果が，積分の計算を支えることになる。

\section*{章末問題}
\begin{enumerate}[label=\textbf{\thechapter.\arabic*},leftmargin=*]
 \item\label{問:極限と漸化式} 次の極限を求めよ。また，$a_1=1$，$a_{n+1}=\sqrt{2+a_n}$ で定まる数列が単調増加かつ上に有界であることを示し，その極限を求めよ。
 \[
  \text{(1)}\quad\lim_{n\to\infty}\frac{2n^2+3}{5n^2-n},
  \qquad
  \text{(2)}\quad\lim_{n\to\infty}\bigl(\sqrt{n^2+2n}-n\bigr).
 \]
 \item\label{問:等比誤差} $\sum_{k=0}^{\infty}\left(-\frac12\right)^k$ の和を求めよ。$k=0$ から $n-1$ までの和で近似するとき，誤差の絶対値を $10^{-3}$ 未満にする最小の自然数 $n$ を求めよ。
 \item\label{問:連続と微分可能} $x\ne0$ で $f(x)=\frac{\sin2x}{x}$，$x=0$ で $f(0)=c$ とする。$f$ が $0$ で連続となる $c$ を求めよ。また，$g(x)=|x-1|$ は $1$ で連続だが微分可能でないことを，差商を使って説明せよ。
 \item\label{問:微分計算} 次の関数を微分せよ。(2) については $x=0$ における接線も求めよ。
 \[
  \text{(1)}\quad\log(1+x^2),\qquad
  \text{(2)}\quad(x^2+1)\e^{-x},\qquad
  \text{(3)}\quad x^{\sin x}\quad(x>0).
 \]
 \item\label{問:平均値評価} 平均値の定理を使い，任意の実数 $a,b$ に対して $|\sin b-\sin a|\le|b-a|$ が成り立つことを示せ。適用する関数と区間を明記せよ。
 \item\label{問:三次関数グラフ} $f(x)=x^3-3x^2$ の増減，極値，凹凸，変曲点を求めよ。さらに，区間 $[-1,3]$ における最大値・最小値と，それらをとる $x$ を求めよ。
 \item\label{問:指数不等式方程式} 微分を用いて，すべての実数 $x$ について $\e^x\ge1+x$ を示し，等号成立条件を求めよ。また，方程式 $\log x+x=2$ は $x>0$ でただ一つの解をもち，その解が $(1,2)$ にあることを示せ。
 \item\label{問:逆媒介陰} 次を求めよ。
 \begin{enumerate}[label=(\arabic*),leftmargin=*]
  \item $f(x)=x^3+x$ の逆関数を $g$ とするとき，$g'(0)$。
  \item $x=t^2$，$y=t^3$ で表される曲線の，$t=2$ における接線。
  \item 曲線 $x^2+xy+y^2=7$ の点 $(1,2)$ における接線。ただし，その点の近くで $y$ は $x$ の微分可能な関数として表されるものとする。
 \end{enumerate}
\end{enumerate}

\smallskip\noindent 解答は\pageref{解答:第2章}ページ。
